% Generated draft: requires mathematical and source review.
\documentclass[11pt]{article}
\usepackage[margin=25mm]{geometry}
\usepackage{fontspec}
\setmainfont{Latin Modern Roman}
\usepackage{amsmath,amssymb,mathtools,unicode-math}
\setmathfont{Latin Modern Math}
\usepackage{enumitem,needspace,xurl,hyperref}
\setlength{\parindent}{0pt}
\setlength{\parskip}{5pt}
\setlength{\emergencystretch}{3em}
\newcommand{\N}{\mathbb{N}}
\newcommand{\Z}{\mathbb{Z}}
\newcommand{\Q}{\mathbb{Q}}
\newcommand{\R}{\mathbb{R}}
\newcommand{\C}{\mathbb{C}}
\newcommand{\F}{\mathbb{F}}
\newcommand{\A}{\mathbb{A}}
\newcommand{\PP}{\mathbb{P}}
\newcommand{\E}{\mathbb{E}}
\newcommand{\eps}{\varepsilon}
\newcommand{\dd}{\,\mathrm{d}}
\begin{document}
\section*{UnsolvedMath ID: 30003224}
\textbf{Source problem code:} catalogue entry
\subsection{Definitions and mathematical statement}
Use the customary holomorphic interpretation of the source statement: let \(n\geq 1\), let \(S\subset\C^n\) be a finite set, and let
\[
f:(\C^n,S)\longrightarrow(\C^{n+1},0)
\]
be a holomorphic multigerm. Two such multigerms are \(\mathcal A\)-equivalent if they differ by biholomorphic changes of source and target coordinates. The germ is \(\mathcal A\)-finite if its \(\mathcal A\)-equivalence class is determined by a finite jet.

Let \(\theta(f)\) be the space of holomorphic vector fields along \(f\), and let \(\theta_{\C^n,S}\) and \(\theta_{\C^{n+1},0}\) denote holomorphic source and target vector fields, respectively, without requiring them to vanish at the marked points. Define the extended normal space and extended codimension by
\[
N\mathcal A_e(f)=
\frac{\theta(f)}
{\{df\circ\xi:\xi\in\theta_{\C^n,S}\}
 +\{\eta\circ f:\eta\in\theta_{\C^{n+1},0}\}},
\qquad
\mathcal A_e\text{-}\operatorname{codim}(f)
=\dim_{\C}N\mathcal A_e(f).
\]
Here the quotient is a quotient of complex vector spaces. For a sufficiently small stable perturbation of a representative of \(f\), let \(X_u\) denote its local image, called a disentanglement. In the setting in which the image Milnor number is defined by this construction, put
\[
\mu_I(f)=\operatorname{rank}_{\Z}\widetilde H_n(X_u;\Z).
\]
When the disentanglement has the homotopy type of a bouquet of \(n\)-spheres, this is the number of those spheres.

\textbf{Proposed statement (Mond's conjecture, as supplied).}
For every such \(\mathcal A\)-finite multigerm in the source's stipulated nice dimensions,
\[
\mathcal A_e\text{-}\operatorname{codim}(f)\leq\mu_I(f).
\]
Moreover, if \(f\) is weighted homogeneous, then
\[
\mathcal A_e\text{-}\operatorname{codim}(f)=\mu_I(f).
\]
Weighted homogeneity customarily means homogeneity under positive weighted coordinate dilations in suitable coordinates. For a single source point, this means that there are positive integer weights \(w_1,\ldots,w_n\) and degrees \(d_1,\ldots,d_{n+1}\) such that
\[
f_j(t^{w_1}x_1,\ldots,t^{w_n}x_n)=t^{d_j}f_j(x)
\qquad(1\leq j\leq n+1).
\]
The supplied record does not specify the compatibility convention for weighted homogeneity across multiple source branches, nor the precise convention or numerical range for nice dimensions. These qualifications remain unresolved; no unrestricted assertion for all \(n\) or all multigerm weight conventions is intended.
\subsection{Short English statement}
For an \(\mathcal A\)-finite holomorphic map multigerm from \(\C^n\) to \(\C^{n+1}\) in the stipulated nice dimensions, the proposed inequality bounds its extended \(\mathcal A\)-codimension by its image Milnor number. Equality is proposed when the germ is weighted homogeneous. The exact dimension range and the multigerm weight convention require source clarification.
\subsection{Sources}
\begin{itemize}
\item Supplied local UnsolvedMath registry record, permanent ID \(30003224\), title: Mond's Conjecture for Finite Map Germs; source problem code \texttt{catalogue entry}; collection: Oberwolfach Reports Open Problems. No individual proposer is supplied. The proposed year is recorded as 2016.
\item Supplied report citation: \emph{Singularities (2017)}, Oberwolfach Reports, DOI \texttt{10.4171/owr/2016/46}, pp.\ 2625--2691. Supplied source URL and literature reference labelled Oberwolfach Report 46: \url{https://doi.org/10.4171/owr/2016/46}. The exact problem page and original passage are not supplied.
\item Literature assessment dated \texttt{2026-08-21}: this is supplied dated curation, not an independently verified assessment. It labels the problem open, says the extracted record does not fully reproduce the exact conjecture, reports that it is known for \(n=1,2\) and remains open for \(n\geq3\), and states that no later primary proof was identified. The supplied background explicitly requests independent verification of time-sensitive status and cited claims.
\item The registry labels the statement \texttt{corrected\_verified} and records a manual semantic, typographic, and LaTeX review claiming that the clean statement is self-contained and preserves the mathematical target. That supplied review label does not resolve the separate warning about incomplete extraction or constitute independent verification here.
\end{itemize}
\end{document}
